.he
.cs 5
.mt 5
.mb 6
.pl 66
.ll 65
.po 10
.hm 2
.fm 3
.sp 2
.sh
2.10  Low-level Organization
.qs
.he PL/I Reference Manual                2.10  Low-level Organization
.tc    2.10  Low-level Organization
.pp 0
.ix low-level organization of PL/I programs
.ix PL/I character set
.ix identifiers
.ix keywords
.ix declared names
.ix operators
.ix declared names
.ix constants
.ix delimiters
.ix comments
The low-level organization of PL/I source text includes a specification of 
the character set and the rules for forming identifiers, both keywords and 
declared names, operators, constants, delimiters, and comments.
.pp
PL/I is a free-format language.  The source program consists of a sequence 
of ASCII characters that make up lines delimited by carriage return characters.
You can enter the source text without regard for column position or specific 
line format.  However, the source text is easier to read and comprehend if you 
follow some basic formatting rules: 
.ix free-format language
.ix source program
.ix ASCII characters
.ix readability of programs
.sp 
.in 5
.ti -2
o Place only one statement on a line.
.sp
.ti -2
o Use indentation to show the nesting level of blocks and DO-groups.
.in 0
.ix DO-group
.ix nested blocks
.bp
.pp
You can create the PL/I source program using any suitable text editor.  
.sp
.sh
Note:\c
.qs
\  all PL/I source programs must have the filetype PLI.
.sp 2
.sh
2.11  The Character Set 
.qs
.he PL/I Reference Manual                     2.11  The Character Set    
.tc    2.11  The Character Set
.pp 
The PL/I character set consists of both upper- and lower-case letters, numeric 
digits, and other symbols.  Table 2-1 shows the symbols recognized by PL/I and 
briefly describes their use.
.sp 2
.ll 60
.ce
.sh
Table 2-1.  PL/I Symbols
.sp
.nf
    Symbol                     Meaning
.in 15
.ti -10
.sp
 =      equal sign (assignment)
.ti -10
 +      plus sign (addition)
.ti -10
 -      minus sign (subtraction)
.ti -10
 *      asterisk (multiplication) 
.ti -10
 /      slash (division) 
.ti -10
 (      left parenthesis (delimiter)
.ti -10
 )      right parenthesis (delimiter)
.ti -10
 ,      comma (separator)
.ti -10
 .      period (name qualifier)
.ti -10
 %      percent symbol (INCLUDE or REPLACE prefix)
.ti -10
 '      apostrophe (string delimiter)
.ti -10
 ;      semicolon (statement terminator)
.ti -10
 :      colon (separator for ENTRY or LABEL constant)
.ti -10
 ^      circumflex (logical Not symbol)
.ti -10
 ~      tilde (alternative Not symbol)
.ti -10
 &      ampersand (logical And symbol)
.ti -10
 |      vertical bar (logical Or symbol)       
.ti -10
 !      exclamation mark (alternative Or symbol)
.ti -10
 \\      backslash (alternative Or symbol)
.ti -10
 >      right angle bracket (greater than)
.ti -10
 <      left angle bracket (less than)
.ti -10
 _      break or underscore (for readablity in identifiers)
.ti -10
 $      dollar sign (valid character in identifiers) 
.ti -10
 ?      question mark (valid character in identifiers)
.fi
.in 0
.ll 65
.bp
.cp 5
.tc    2.12  Identifiers
.sh
2.12  Identifiers
.qs
.he PL/I Reference Manual                           2.12  Identifiers
.pp
An identifier is a string of from one to thirty-one characters that are either 
letters, digits, or the underscore.  The first character must be a letter.  
PL/I always represents letters internally in upper-case.  Therefore, two 
identifiers that differ only in case represent the same identifier.
.ix identifier, maximum length
.ix identifier, formation of
.pp
PL/I allows the question mark character to be embedded in identifiers to allow 
access to external system entry points.   
.sp
.sh
Note:\c
.qs
\  you should avoid embedded question marks to maintain upward 
compatibility with full language implementation.
.ix external system entry points
.pp
Every identifier in the source text of a PL/I program must be either a keyword 
or a declared name.  Keywords are those identifiers that have a special 
meaning in PL/I when used in a specific context.  Examples of keywords are the 
names of built-in functions, statements, and data attributes.  The following 
is a list of all the keywords.  The \c
.ul
PL/I Language Command Summary\c
.qu
\ contains a complete list of keywords with brief explanations.  
.ix PL/I keyword
.ix declared name
.sp 2
.nf
     A              ABS            ACOS              ADDR                
     ALIGNED        ALLOCATE       ASCII             ASIN
     ATAN           ATAND          AUTO              AUTOMATIC
     B              B1             B2                B3
     B4             BASED          BEGIN             BIN
     BINARY         BIT            BOOL              BUILT-IN
     BY             CALL           CEIL              CHAR
     CHARACTER      CLOSE          COLLATE           COLUMN
     COS            COSD           DCL               DEC               
     DECIMAL        DECLARE        DIM               DIMENSION         
     DIRECT         DIVIDE         DO                E              
     EDIT           ELSE           END               ENDFILE        
     ENDPAGE        ENTRY          ENV               ENVIRONMENT    
     ERROR          EXP            EXT               EXTERNAL
     F              FILE           FIXED             FIXEDOVERFLOW
     FLOAT          FLOOR          FOFL              FORMAT
     FREE           FROM           GET               GO TO             
     GOTO           HBOUND         IF                INCLUDE           
     INDEX          INIT           INITIAL           INTO
     KEY            KEYED          KEYFROM           KEYTO
     LABEL          LBOUND         LENGTH            LINE
     LINENO         LINESIZE       LIST              LOG
     LOG2           LOG10          MAIN              MAX
     MIN            MOD            NULL              OFL
     ON             ONCODE         ONFILE            ONKEY               
     OPEN           OPTIONS        OUTPUT            OVERFLOW            
     PAGENO         PAGESIZE       POINTER           PRINT
     PROC           PROCEDURE      PTR               PUT
.bp
     R              RANK           READ              RECORD
     RECURSIVE      REPEAT         REPLACE           RETURN
     RETURNS        REVERT         ROUND             SEQUENTIAL
     SET            SIGN           SIGNAL            SIN
     SIND           SINH           SKIP              SQRT
     STACK          STATIC         STOP              STREAM
   ~  SUBSTR         SYSIN          SYSPRINT          TAB
     TAN            TAND           TANH              THEN
     TITLE          TO             TRANSLATE         TRUNC
     UNDEFINEDFILE  UNDF           UNDERFLOW         UFL
     UNSPEC         UPDATE         VAR               VARIABLE
     VARYING        VERIFY         WHILE             WRITE
     X              ZERODIVIDE     
.sp 
.fi
.pp
Declared names are identifiers whose use or meaning you define in a DECLARE 
statement (Section 3.6).  A keyword can appear in a declaration as a 
user-defined identifier.  The meaning of the identifier depends on how and 
where it appears.  PL/I determines the meaning in context.  For example, INDEX 
is a keyword because it is the name of a PL/I built-in function.  However, in 
the context of the declaration,
.ix context of a declaration
.sp
.in 5
declare index fixed binary;
.sp
.in 0
index is a declared name and not a keyword.
.sp 2
.sh
2.13  Constants
.qs
.he PL/I Reference Manual                             2.13  Constants
.tc    2.13  Constants
.pp
.ix constants
Constants are text items that have a fixed literal meaning which can not 
change when the program runs.  In PL/I, the basic constants are the following:
.sp
.nf
.in 3 
o arithmetic  (Example: 3674.799)
o character string  (Example: 'Ada Lovelace')
o bit string  (Example: '00010110'B)
.fi
.in 0
.ix constant, arithmetic
.ix constant, character string
.ix constant, bit
.sp 2
.sh
2.14  Delimiters and Separators
.qs
.he PL/I Reference Manual             2.14  Delimiters and Separators
.tc    2.14  Delimiters and Separators
.pp
Separate items, such as identifiers, must be distinguishable.  PL/I recognizes 
certain characters as delimiters and separators.  
.ix delimiter 
.ix separator
.pp
Usually, delimiters enclose one or more text items while separators mark the 
end of one item and the beginning of another.  In PL/I, each identifier and
arithmetic constant must be preceded and followed by one or more delimiters or 
separators.  Delimiters can be either spaces, operators, or certain special 
characters.
.bp
.sh
2.14.1  Spaces
.qs
.tc         2.14.1  Spaces
.ix spaces
.ix blanks
.ix tab characters, CRTL-I
.pp
In PL/I, a space can be either a blank, or a tab character (CTRL-I).  PL/I 
ignores any carriage return, line-feed, or carriage return, line-feed sequence 
that is embedded in a string constant.  For example, the assignment statement,
.ix carriage return
.ix line-feed
.ix carriage return line-feed pair
.sp
string = 'WHEN YOU HAVE A VERY LONG STRING LIKE THIS, PL/I
ALLOWS YOU TO PUT SOME OF IT ON ANOTHER LINE';
.sp
.in 0
assigns the specified character string to the variable string.  Any blanks or 
tabs that follow ALLOWS or precede YOU TO PUT are included in the string. 
.sp 2
.sh
2.14.2  Operators
.qs
.tc         2.14.2  Operators
.ix operator, symbol for a mathematical of logical operation 
.pp
An operator is a symbol for a mathematical or logical operation.  There are 
four types of operators in PL/I as shown in Table 2-2.
.sp
.sh
Note:\c
.qs
\  operators that consist of two characters, such as >=, are called composite 
operators and must not be separated by blanks or tabs.
.ix composite operator
.bp
.sh
                  Table 2-2.  PL/I Operators
.qs
.sp
.nf
              Symbol              Meaning    
.fi
.in 0
.sp
.ll 60
.in 25
Arithmetic Operators
.sp
.in 25
.ti -6
+     addition or prefix plus
.ti -6
-     subtraction or prefix minus
.ti -6
*     multiplication
.ti -6
/     division
.ti -6
**    exponentiation
.sp 2
.in 25
Comparison Operators
.sp
.in 25
.ti -6
>     greater than
.ti -12
^> or ~>    not greater than
.ti -6
>=    greater than or equal to
.ti -6
=     equal to
.ti -12
^= or ~=    not equal to
.ti -6
<=    less than or equal to
.ti -6
<     less than
.ti -12
^< or ~<    not less than
.sp 2
.in 25
Bit-string Operators 
.sp
.in 25
.ti -11
^ or ~     Logical Not
.ti -6
&     Logical And
.ti -13
| or ! or \\  Logical Or
.sp 2
.in 0
.in 25
The String Operator
.sp
.in 25
.ti -16
|| or !! or //  concatenate
.ix concatenation of strings
.ix concatenation operator
.sp 3
.ll 65
.in 0
.sh
2.14.3  Special Characters 
.qs
.tc         2.14.3  Special Characters 
.pp
Table 2-3 shows the special characters that can also function as delimiters
or separators in PL/I.  Subsequent sections of the manual contain examples
of their use.  
.ix special characters
.bp
.in 5
.sh
Table 2-3.  Special Character Delimiters and Separators
.qs
.sp
.in 0
.ll 60
.nf
   Character                       Function
.fi
.sp
.in 19
.ti -12 
:           A colon follows ENTRY and LABEL constants.
.sp
.ti -12
;           A semicolon terminates statements.
.sp
.ti -12
,           A comma separates elements of a list.
.sp
.ti -12
.cc *
.           A period separates items in a qualified name.
*cc .
.sp
.ti -12
'           A single apostrophe is a delimiter for the specification of 
character and bit-string constants.
.sp
.ti -12
->          The arrow is a composite operator consisting of the minus sign and 
the right angle bracket.  It is a separator in a pointer qualified reference.
.sp
.ti -12
=           An equal sign is a separator in an assignment statement.
.sp 
.ti -12
(           Left parenthesis.
.sp
.ti -12
)           Right parenthesis.  A left parenthesis together with a right 
parenthesis is used to enclose lists and extents, define the 
order of evaluation of expressions, and separate keywords from 
statements and option names.
.in 0
.ll 65
.sp 3
.sh
2.14.4  Comments
.qs
.tc         2.14.4  Comments
.pp
Comments provide documentary text in a PL/I source program.  The compiler 
ignores comments, so you can place them wherever a delimiter is appropriate.  
Precede a comment by the composite pair /* and end the comment by the reverse 
composite pair */.  For example,
.ix comments
.ix documentary text
.ix source program
.sp
.in 5
.nf
   .                                            
   .                                            
   .                                            
get list(name);  /* read the name */            
   .                                            
   .                                            
   .                                            
.in 0
.fi
.bp
.sh
2.15  Preprocessor Statements
.he PL/I Reference Manual               2.15  Preprocessor Statements
.tc    2.15  Preprocessor Statements
.pp
PL/I allows modification of the source program or inclusion of external source 
files at compile time through the use of preprocessor statements.  
Preprocessor statements are identified by a leading % symbol before the 
keyword:
.ix source program
.ix preprocessor statements
.ix statements, preprocessor
.sp
.in 5
INCLUDE   or   REPLACE
.in 0
.sp 2
.sh
2.15.1  The %INCLUDE Statement
.qs
.tc         2.15.1  The %INCLUDE Statement
.pp
The %INCLUDE statement copies PL/I source text from an external file at 
compile time.  The statement is useful for filling in a structure declaration 
or format list.  The %INCLUDE statement has the form: 
.ix %INCLUDE statement
.ix statements, %INCLUDE
.ix source text
.sp
.ti 5
%INCLUDE 'filespec';
.sp
where filespec designates the file to copy into the source program.  Filespec 
must be a standard file specification, [d:]filename[.typ], and must
be enclosed in single apostrophes.  If there is no drive specification, PL/I 
assumes the drive containing the source program.  When the compiler 
encounters the %INCLUDE statement in the source file, it begins reading the 
file specified by %INCLUDE.  When the compiler reaches the end of the 
%INCLUDE file, it resumes reading the original source file.
.ix file specification
.pp
The following code sequence is an example of the %INCLUDE statement: 
.sp 2
.in 5
.nf
f:                                       
  procedure;                             
  declare a fixed binary;                
  %include 'struc.lib';                  
  declare c float;                       
    .                                    
    .                                    
    .                                    
end f;                                   
.sp 2
.in 0
.fi
The compiler includes the source text from the file struc.lib at the point 
of the %INCLUDE statement.
.sp
.sh
Note:\c
.qs
\  PL/I does not allow nested %INCLUDE statements. 
.ix nested %INCLUDE statements
.sp 2
.mb 4
.fm 1
.sh
2.15.2  The %REPLACE Statement
.qs
.tc         2.15.2  The %REPLACE Statement
.ix %REPLACE statement
.ix statements, %REPLACE
.pp
The %REPLACE statement allows you to program with named constants.  The 
%REPLACE statement has the form: 
.sp
.ti 5
%REPLACE identifier BY constant;
.sp
The compiler replaces every occurrence of the given identifier in the source 
text with the specified constant.  The constant can be a signed or unsigned 
arithmetic constant, a bit string, or a character string.  You can write 
multiple %REPLACE statements as a single %REPLACE statement, with the 
elements separated by commas.
.ix identifier
.ix multiple %REPLACE statements
.pp
.mb 6
.fm 3
For example, the statement
.sp
.in 5
%replace true by '1'b;
.in 0
.sp
replaces all occurrences of true by the bit string constant '1'b, so that the 
compiler interprets the statement,
.sp
.ti 5
do while(true);
.sp
as: 
.sp
.ti 5
do while('1'b);
.pp
PL/I requires that all %REPLACE statements occur at the outer block level 
before any nested inner blocks.  
.sp
.sh
Note:\c
.qs
\  to facilitate program maintenance and debugging, you should write all 
%REPLACE statements directly following the procedure heading. 
.ix nested blocks
.sp 2
.ce
End of Section 2
.nx three.tex

